\documentclass{article}
\usepackage{amsmath, amssymb, amsthm}
\usepackage{hyperref}
\usepackage{geometry}
\geometry{margin=1in}

\title{Erd\H{o}s Problem \#528}
\date{Last updated: 2025-08-31}
\author{Source: erdosproblems.com}

\begin{document}
\maketitle

\noindent\textbf{Status:} OPEN \quad \textbf{No prize}

\noindent\textbf{Tags:} geometry

\medskip
\noindent\textbf{Problem Statement:}

Let $f(n,k)$ count the number of self-avoiding walks of $n$ steps (beginning at the origin) in $\mathbb{Z}^k$ (i.e. those walks which do not intersect themselves). Determine\[C_k=\lim_{n\to\infty}f(n,k)^{1/n}.\]

\bigskip
\noindent\textbf{Remarks:}

The constant $C_k$ is sometimes known as the connective constant. Hammersley and Morton \cite{HM54} showed that this limit exists, and it is trivial that $k\leq C_k\leq 2k-1$.

Kesten \cite{Ke63} proved that $C_k=2k-1-1/2k+O(1/k^2)$, and more precise asymptotics are given by Clisby, Liang, and Slade \cite{CLS07}.

Conway and Guttmann \cite{CG93} showed that $C_2\geq 2.62$ and Alm \cite{Al93} showed that $C_2\leq 2.696$. Jacobsen, Scullard, and Guttmann \cite{JSG16} have computed the first few decimal places of $C_2$, showing that\[C_2 = 2.6381585303279\cdots.\]See also [529].

\bigskip
\noindent\textbf{References:}

[Al93] Alm, Sven Erick, Upper bounds for the connective constant of self-avoiding
walks. Combin. Probab. Comput. (1993), 115-136.
 
 [CG93] Conway, A. R. and Guttmann, A. J., Lower bound on the connective constant for square lattice
self-avoiding walks. J. Phys. A (1993), 3719-3724.
 
 [CLS07] Clisby, Nathan and Liang, Richard and Slade, Gordon, Self-avoiding walk enumeration via the lace expansion. J. Phys. A (2007), 10973-11017.
 
 [HM54] Hammersley, J. M. and Morton, K. W., Poor man's Monte Carlo. J. Roy. Statist. Soc. Ser. B (1954), 23-38; discussion 61-75.
 
 [JSG16] Jacobsen, Jesper Lykke and Scullard, Christian R. and
Guttmann, Anthony J., On the growth constant for square-lattice self-avoiding walks. J. Phys. A (2016), 494004, 18.
 
 [Ke63] Kesten, Harry, On the number of self-avoiding walks. J. Mathematical Phys. (1963), 960-969.

\bigskip
\noindent\small{Source: \url{https://www.erdosproblems.com/528}}
\end{document}
